\documentclass[10pt]{article}
\usepackage{fontspec}
\setmainfont{lmroman10-regular.otf}[BoldFont=lmroman10-bold.otf,ItalicFont=lmroman10-italic.otf,BoldItalicFont=lmroman10-bolditalic.otf]
\setmonofont{lmmono10-regular.otf}
\usepackage[a4paper,margin=16mm,top=14mm,bottom=16mm]{geometry}
\usepackage{multicol,graphicx,xcolor,array,fancyhdr}
\usepackage[hidelinks]{hyperref}
\setlength{\columnsep}{6mm}
\setlength{\parindent}{1em}
\setlength{\parskip}{2pt}
\newcommand{\unk}[1]{\textcolor{red}{\texttt{\scriptsize[U+#1]}}}
\pagestyle{fancy}\fancyhf{}\renewcommand{\headrulewidth}{0.4pt}
\fancyhead[C]{\small 5550 \quad \textbf{Part II.}}
\begin{document}
% block 1 Title bbox=[0.4004, 0.0453, 0.627, 0.0648]
\section*{\centering \textbf{OPTICS.}}
\begin{multicols}{2}\raggedcolumns
% block 3 Caption bbox=[0.0596, 0.0688, 0.125, 0.0984]
\begin{center}\small\itshape or Refletions.\end{center}
% block 4 Text bbox=[0.1221, 0.0664, 0.5127, 0.318]
tre is C, and focus of parallel rays is at P; and let the imaginary focus of the incident rays be as F, a point between P and B, and he DR be the refleth ray. From C and R draw the lines CH, RH, the one palling through D, the other parallel to FE. Then will the angle RHD be equal to HDD the angle of incidence, as alterant to is; and therefore equal to HDD, the angle of reflections is reflected the triangle HDR is indicates, and consequently DR is equal to RH. Now the lines FD and RH being parallel, the complex DDC and RHD are similar; and therefore RH/ or RD, is to PD as CR on CF: but BD vanthing, RD and FD coincide with RB and IB, whether RF is to TR as CR in CF; that is, the distance of the focus from the surface is to the dilance of the imaginary focus of the incident rays, at the distance of the focus from the centre is to the dilance of the imaginary focus of the incident rays from the fame. \_@E.D.\_

% block 5 Text bbox=[0.123, 0.3187, 0.5117, 0.357]
Cass. Of diverging rays falling upon a concave surface, and corresponding from a point between the fo-cus of parallel rays and the surface.

% block 6 Text bbox=[0.124, 0.3578, 0.5137, 0.5156]
Dass. Let ED and FB represent two diverging rays flowing from the point 3 as a radiant, which was the imaginary focus of the incident rays in the forgoing cafe. Then an SD as in the cafe selected with anDR, FD will be reflected into DR(, (for the region unemtioned in Cafc the feecond), lo that the refleted ray will proceed as from the point 3, but it cane demonflatred in the cafe manufacturing forceing, that RB is to FB as CR to CF; that is, the distance of the fo-cus from the surface is to that of the radiant from the fame, as the distance of the former from the centre is to that of the latter from the fame. \_@E.D.\_

% block 7 Text bbox=[0.1201, 0.5172, 0.5146, 0.5555]
Cass 7. Of converging rays falling upon a convex surface, and including towards a point beyond the centro.

% block 8 Text bbox=[0.1191, 0.5563, 0.5156, 0.7922]
Dass. Let AB, ED, (fig. 8) be the incident rays tending to F, a point beyond the centre C, and let DR be the refleted ray of the incident one ED. Then because the incident ray ED tend to a point beyond the centre, the refleted ray DR will proceed as from one on the contrary to, flopode 2 as asrelated above under Drop, VII. Through D draw the perpendicular CD, and produce it to H. Thus will EDH and HDD be the angles of incidence and reflelections, while being equal, their vertical ones of DF and CDR will be fo to cases confiantly the vertex of the triangle FDR and bifled by the line CD+ wherefore, RD is to DF, or (3 Elem d o) BD vanthing, RB is to DF as RC to CF; that is, the dilatance of the focus of the refleted rays is to that of the imaginary focus of the incident rays, as the dilatance the lower from the center is to the dilatance of the latter from the fame. \_@E.D.\_

% block 9 Text bbox=[0.1338, 0.7937, 0.5166, 0.8336]
Cass 8. Of diverging rays falling upon a concave surface, and proceeding from a point beyond the centro.

% block 10 Text bbox=[0.1348, 0.8328, 0.5176, 0.9289]
Dass. Let FB, ED, be the incident rays having their radiant in F, the imaginary focus of the incident rays in the foreground cafe. Then an ED was in that cafe reflected into DR(, FD will now be reflected into DR; lo that R will be the focus of the refleted rays. But it was demontrated in the foregoing cafe, that RB is to FB as RC to CF; that is, the dilatance of

% block 11 Text bbox=[0.5195, 0.0656, 0.9775, 0.1336]
the focus of the refleted rays from the furface is to the dilatance of the radiant from the fame, as the diflance of the focus of the refleted rays from the centre is to the dilatance of the radiant from thence. \_@E.D.\_

% block 12 Text bbox=[0.5205, 0.1344, 0.9062, 0.2]
The two remaining cafes may be confirmed as the converfe of thote under Prop. II. (p. 548-), because the incident rays in there are the reflected ones in them; or they may be demothered in the fame manner with the forgoing, as follows.

% block 13 Text bbox=[0.5215, 0.2008, 0.9053, 0.2414]
Cass 9. Converging rays falling upon a convex furface, and tending to the surface of parallel rays, become parallel after reflection.

% block 14 Text bbox=[0.5225, 0.2406, 0.9072, 0.4773]
Dass. Let ED, RB, (fig. 8) reported two convergency rays inside to the convex fur- End, and tending towards F, which we will now loppole to be the focus of parallel rays; and he DR be the refleted ray and Cash center of coronavity of the reflefling surfaces. Through C draw the line CD, and produces to FR, drawing ARH parallel to ED prod-cut to F. Now it has been demottemical Cade ? when the incident rays are loppodof to tend to the point Y, that RB is to FB as RC to CF; but F in the Cafe being loppold to be the focus of parallel rays, it is the middle point between C and B (by proposition all), and therefore FB and FC are equal and confiantly the two outer corners in the proportion, viz. RB and RC, mitb fe toto ; which can only be upon a loppodiation that R is at an infinite diffuse from B; that is, that the refleted ray DR and DR be parallel. \_@E.D.\_

% block 15 Text bbox=[0.5234, 0.4781, 0.9082, 0.5164]
Cass to Diverging rays falling upon a concave fur-ness and preceding from the focus of parallel rays, become parallel after reflection.

% block 16 Text bbox=[0.5244, 0.518, 0.9092, 0.6344]
Dass. Let AD, EB, BB (fig. 6) be two diverging rays incident upon the concave fur BD, as loppofed in Cafe 4; where it was demottemical that FB is to RB as CF to OR, then in the prefect each RB and CR are equal, because R is loppold to be the focus of parallel rays; therefore FB and FC are fo too. Which cannot be under F be taken as the nature difame from B; that is, and the refleted rays BF and DF be parallel. \_@E.D.\_

% block 17 Text bbox=[0.5254, 0.6359, 0.9111, 0.9195]
It is here otherable, that in the cafe of diverging rays falling upon a convex furace, (fe fig. 7) the further the point It is taken from B, not never the point F, the focus of the refleted rays, approaches to B, while the radiant R remains the fame. For is in evident from the curvature of events, thus the ripet D (fig. 9) may be te ken fo far from B, that the refleted ray DE final proceed as from F, G, H, or even from B, or from points between B and R; and the further is in taken from B, the father the point from which it prococis, appochoz toward R is will easily appear if we doue forced incident rays with that respective reflected ones, in fashion than that angles of reflection may be all equal to their refleting angles of incidentes as at those in the figure. The fact is applicable to any of the other cafes of diverging or converging rays incident upon a spherical surface. That is the reason, that when rays are considered as esflected from a spherical surface, the dilatance of the oblique rays from the perpendicular one is taken to [all, that it may be loppold to vanish.

% block 18 Text bbox=[0.8672, 0.9203, 0.9102, 0.9375]
From

\end{multicols}
\vfill\par\noindent\rule{\linewidth}{0.4pt}\par
{\scriptsize\color{gray}\noindent Source: Encyclopaedia Britannica scan, text from Parse 3.5 without correction. page\_id: britannica:309f765b1ecc83e0a7f7762e44b4a927a7555443:train:851 \quad row index: 851 \quad image sha256: 6e2a0b9993242da4f83259432110d6929e728a4c0e58a9e79a16ff80a6fd1861 \quad parse model: nvidia/nemotron-3.5-parse \quad parse finish: stop \quad scan: \href{../scans/page-0851-6e2a0b999324.jpg}{page-0851-6e2a0b999324.jpg}\par}
\end{document}
